\documentclass[10pt,a4paper]{amsart}
\usepackage[margin=19mm]{geometry}
\usepackage{amsmath,amssymb,mathtools}
\usepackage[scr=rsfs,cal=euler]{mathalfa}
\usepackage{xfrac}
\usepackage[unicode,hidelinks,bookmarks=false]{hyperref}
\newcommand{\ie}{i.e.\ }
\newcommand{\cf}{cf.\ }
\newcommand{\resp}{respectively\ }
\newcommand{\Subsection}{\subsection}
\providecommand{\nobreakdash}{\nobreak\hbox{-}}
\setlength{\emergencystretch}{2em}
\multlinegap0pt
\usepackage{bm}
\usepackage{bbm}
\usepackage{dsfont}
\usepackage{xspace}
\usepackage{cancel}
\usepackage{mathtools}
\usepackage{amsthm}
\makeatletter
\@ifundefined{lemma}{\newtheorem{lemma}{Lemma}}{}
\makeatother
\newcommand{\EE}{\mathbb E}
\newcommand{\PP}{\mathbb P}
\newcommand{\Var}{\operatorname{Var}}
\title[Estimating Community Boundaries in Geometric Random Graphs]{Estimating Community Boundaries in Geometric Random Graphs}
\author{Taha Ameen, Neeladri Maitra}
\date{Source: arXiv:2608.12054v1 (CC BY 4.0)}
\begin{document}
\maketitle

\noindent\emph{Source and licence.} Estimating Community Boundaries in Geometric Random Graphs. Version arXiv:2608.12054v1 (CC BY 4.0), distributed under the Creative Commons Attribution 4.0 International licence, as recorded in the arXiv licence metadata for this exact version and shown on the abstract page https://arxiv.org/abs/2608.12054v1. Full rights notice: licenses/arxiv-2608.12054-CC-BY-4.0.txt.\par\smallskip

\noindent\textit{Editorial scope.} Counted unit: the Lemma whose source label is \texttt{lem:localization} in the article (source0.tex lines 880--908, source label \texttt{lem:localization}), delivered together with the article's own complete original proof of it (source0.tex lines 911--1041, one \texttt{proof} environment of 131 lines). No mathematics is written, repaired or re-derived: statement and proof are the article's own text line for line. Required context retained uncounted: lem:local-endpoint (lines 736--757). Every definition, display and label of those lines is retained unchanged. Omitted from this package and not redistributed: 1--735, 758--879, 909--910, 1042--1504. Nothing inside a retained excerpt is omitted or altered: every retained body is delivered exactly as the article's lines are written, so no omission receipt of any kind accompanies this package. Citations: the retained excerpts carry no citation command at all, so no bibliography entry of the article is transcribed and no reference is redistributed. Reference adaptation: none was needed and none was made; every reference inside the delivered unit resolves inside it, so no unresolved reference or citation is produced. Printed slips kept literal: no printed word, symbol, hypothesis or number of the counted statement or of its proof was altered. Layout and class: the article's own class is \texttt{ejpecp} with packages including setspace, enumitem, xcolor, cleveref, natbib, subcaption; the wrapper is the lane's pinned \texttt{amsart} editorial wrapper at 10pt on A4, which restates the theorem-like environments the retained text needs and adds the article's own macro definitions that the retained text needs (\texttt{\textbackslash EE}, \texttt{\textbackslash PP}, \texttt{\textbackslash Var}), with extra wrapper packages (bm, bbm, dsfont, xspace, cancel, mathtools). The delivered numbering is the unit's own. Assets: no retained span contains a figure environment, a graphics-inclusion command, a PDF-page inclusion, a table or a TikZ picture; no image file is copied into this package, the package declares no asset dependency and no image file is redistributed with it. Licence: version arXiv:2608.12054v1 of this article is distributed under the Creative Commons Attribution 4.0 International licence, as recorded in the arXiv licence metadata for this exact version and shown on the abstract page https://arxiv.org/abs/2608.12054v1. A full rights notice is delivered at licenses/arxiv-2608.12054-CC-BY-4.0.txt. Characters: every retained excerpt is pure ASCII, so the delivered engine, XeLaTeX with its default Latin Modern text font, reports no missing character.
\begin{lemma}[Local endpoint probability] \label{lem:local-endpoint}
    Fix \(z\in (0,1)\), and assume \(0\le s<t\le 1/4\).  There exist constants \(c_0>0\), and \(0<c_1<c_2<\infty\) depending only on $z$ such that, whenever
    \(
        0<\delta\le c_0t \, ,
    \)
    one has
    \begin{equation} \label{eq-local-ep}
        c_1\delta(t^2-s^2)
        \le
        \PP\left(
        U^{(1)}<z<V^{(1)}<z+\delta,\;
        s<\|U-V\|\le t
        \right)
        \le
        c_2 \, \delta(t^2-s^2) \, ,
    \end{equation}
    uniformly in \(s,t,\delta\).  The analogous estimate with
    \(
        z-\delta<U^{(1)} < z < V^{(1)}
    \)
    also holds.
\end{lemma}

\begin{lemma}[Endpoint localization] \label{lem:localization}
Assume that \(0\le s_N<t_N\le1/4\) for all sufficiently large \(N\), and fix \(z\in(0,1)\). Let \(a_N = t_N^2 - s_N^2 \), \(b_N = t_N^3 - s_N^3 \) and 
\(
    \Delta_N
    =
    \frac1N+\frac{1}{N^2a_N}.
\) 
If
\(
    N^2 b_N\to\infty \, ,
\)
then, for every fixed \(M\ge1\),
\[
    \PP\left(
        K_N^+(z,M\Delta_N)=0
    \right)
    \, \le \,
    \frac{C_z}{M}+o(1) \, ,
    ~~~~
    \mbox{and}
    ~~~~
    \PP\left(
        K_N^-(z,M\Delta_N)=0
    \right)
    \, \le \, 
    \frac{C_z}{M}+o(1) \, ,
\]
where \(C_z<\infty\) depends only on the fixed point \(z\).
\end{lemma}

\begin{proof}
We prove the statement for \(K_N^+\). Throughout, denote
\(
    s=s_N, \; 
    t=t_N, \;
    a=t^2-s^2, \;
    b=t^3-s^3,
\)
and set
\[
    \delta=M\Delta_N
    =
    M\left(
        \frac1N+\frac{1}{N^2a}
    \right).
\]
Since
\(
    b/a
    =
    {(t^2+ts+s^2)}/{(t+s)}
    \le 3t,
\)
the assumption \(N^2b\to\infty\) implies \(N^2ta\to\infty\).
Moreover, \(b\le t^3\), so \(N^2t^3\to\infty\). Hence
\(Nt^{3/2}\to\infty\), and since \(t\le1/4<1\), it follows that
\(Nt\ge Nt^{3/2}\to\infty\).
%
%
Therefore, for each fixed \(M\),
\[
    \frac{\delta}{t}
    =
    M\left(
        \frac{1}{Nt}
        +
        \frac{1}{N^2at}
    \right)
    \longrightarrow0.
\]
Thus, for all sufficiently large \(N\), we have
\(\delta\le c_0(z)t\) and may apply
Lemma~\ref{lem:local-endpoint}.

Let
\(
    q \coloneqq \EE \big[ E^+_{12} \big] .
\)
The lower bound in Lemma~\ref{lem:local-endpoint} gives
\(
    q\ge c_z\delta a
\)
for some \(c_z>0\). Hence, writing
\[
    \mu=\EE K_N^+(z,\delta)=N(N-1)q,
\]
we have
\(
    \mu\ge c_zN^2\delta a
\)
for all sufficiently large \(N\).

We next bound the variance. Indicators involving disjoint vertex sets are independent. 
If two distinct indicators share a vertex that one event requires to
lie to the left of \(z\) and the other requires to lie to the right of
\(z\), then the two events are disjoint, so the indicators cannot both equal one and their covariance is non-positive.
%
It remains to consider indicators sharing a vertex in the same role. For example, define
\[
    q_L(x) \coloneqq \EE\left[E^+_{12}\mid X_1=x\right].
\]
Conditional on \(X_1\), the indicators \(E^+_{12}\) and \(E^+_{13}\) are independent. Moreover, for every fixed \(x\), an admissible right endpoint must lie in a subset of the annulus centered at \(x\), whose area is \(\pi a\). Thus \(q_L(x)\le\pi a\), and therefore
\[
    \EE\left[E^+_{12}E^+_{13}\right]
    =
    \EE\left[q_L(X_1)^2\right] 
    \le
    \pi a\,\EE \left[ q_L(X_1) \right]
    =
    \pi a q \, .
\]
The same argument applies to two indicators sharing their right endpoint. Since there are at most \(N(N-1)(N-2)\) covariance terms of each type,
\[
    \Var\left(K_N^+(z,\delta)\right)
    \le
    \mu+2\pi N(N-1)(N-2)aq 
    =
    \mu\left(1+2\pi(N-2)a\right) \, .
\]
Consequently,
\[
    \frac{\Var(K_N^+(z,\delta))}{\mu^2}
    \le
    C_z\left(
        \frac{1}{N^2\delta a}
        +
        \frac{1}{N\delta}
    \right).
\]
Finally, note that
\(
    N^2\delta a=M(Na+1)\ge M,
\)
and
\(
    N\delta
    =
    M\left(1+\frac{1}{Na}\right)
    \ge M.
\)
Therefore, Chebyshev's inequality gives
\[
    \PP\left(K_N^+(z,\delta)=0\right)
    \le
    {
    \PP\left( |K_N^+(z,\delta)-\mu|>\mu/2 \right) 
    }
    \le 
    \frac{4\Var(K_N^+(z,\delta))}{\mu^2}
    \le
    \frac{C_z}{M},
\]
after enlarging \(C_z\) if necessary.


%
Taking
\(\delta=M\Delta_N\) proves the claimed bound. The proof for \(K_N^-\)
is the same, using the left-endpoint estimate in
Lemma~\ref{lem:local-endpoint} and enlarging \(C_z\) if necessary.
\end{proof}

\end{document}
